% ---------------------------------------------------------------------
%  6. The hard frontier: sampling scale + the hollow-bank boundary
% ---------------------------------------------------------------------
\section{How far the frontier reaches}
\label{sec:frontier}

The integrated agent solves what is reliably solvable under caps. But ``reliably''
is not the whole story: across configurations the system has produced verified,
model-generated proofs for \textbf{6 of 9} development problems and \textbf{10 of
16} sample problems --- reaching a genuine IMO multi-theorem problem. This section
is honest about both how it gets there and where it stops.

\subsection{Hard wins come from test-time sampling scale}
\texttt{p09\_imo1964} is an IMO problem with a multi-theorem statement, and the
system closed it substantively on a directed run: model-generated,
comparator-clean, no hardcoding. The mechanism is test-time scale, not a cleverer
prompt. Each sub-theorem is sampled independently many times (in that campaign,
roughly one of ten samples passed per sub-theorem) and the accepted pieces are
combined into one file.
Sampling each theorem \emph{separately} breaks the conjunction that dooms a single
long attempt: needing every sub-goal clean in the same generation is
exponentially unlikely, but needing each clean in \emph{some} generation is
cheap. This mirrors the transferable lesson from the wider literature --- with no
fine-tuning, the lever is search and verification at inference time.

It also explains an apparent paradox --- \emph{if the ladder solved \texttt{p09}
once, why not every time?} These hard wins are \emph{probabilistic}: a one-in-ten
per-sample rate makes the outcome depend on how much sampling the budget buys, and
the integrated agent spreads its budget across every rung rather than pouring it
into heavy per-theorem sampling. So \texttt{p09} is ``substantively solvable, not
yet reliable,'' and the fix is a scheduling choice (\S\ref{sec:discussion}), not a
missing capability.

\subsection{The boundary: verified progress is not yet a proof}
For the remaining hard problems I built the most literature-grounded idea
available under the no-retrieval rule: instead of proving the whole theorem, mint
small \emph{verified} lemmas and reuse them, accumulating a bank of checked
partial progress. The bank works as a bank: on \texttt{rmo\_2000\_2} it saved
several lemmas, reused most of them, and recorded genuinely \emph{decisive}
reuses. What it does not do --- consistently --- is assemble those pieces into a
proof of the locked goal. An overnight campaign on the hardest six problems
finished with the bank filling, greedy assembly firing, and the target theorem
still open.

I call this the \textbf{hollow-bank boundary}: the system mints and reuses the
verified ``screws'' but cannot always assemble the ``furniture.'' It names the
failure precisely --- not proof \emph{generation}, which works, but goal-directed
\emph{assembly} --- with verified evidence rather than a vague capability ceiling,
and it is a boundary, not a wall: the same cheap ladder that stops here is the one
that reaches \texttt{p09}. (For completeness, \texttt{rmo\_2000\_2} is
\emph{reachable} by a hand-built cast-and-squeeze proof, deliberately kept out of
the agent since the rules forbid per-problem hardcoding; the open question is
whether the models rediscover the idiom through the generic sampling interface.)
